% =============================================================================
% Module 6: Solutions to Exercises
% =============================================================================
% This file is conditionally included based on \ifsolutions
% =============================================================================

\section{Solutions to Exercises}
\label{sec:fermat_solutions}

\noindent
Full worked solutions are also available in the Mathematica script
\solnlink{06_Fermat_Time_Delays/problems\_06.wl}{problems\_06.wl},
which verifies each answer symbolically and numerically.

% ----- Exercise 6.1 -----
\subsection*{Exercise~\ref{ex:fermat_lens_eq}: Deriving the Lens Equation from Fermat's Principle}

\textbf{(a)} The Fermat potential is
$\tau(\vectheta, \vecbeta) = \half|\vectheta - \vecbeta|^2
- \psiL(\vectheta)$.  Taking the gradient with respect to $\vectheta$:
\begin{equation*}
    \nabla_{\vectheta}\,\tau
    = \nabla_{\vectheta}\!\left[\half|\vectheta - \vecbeta|^2\right]
    - \nabla_{\vectheta}\,\psiL(\vectheta)
    = (\vectheta - \vecbeta) - \nabla\psiL(\vectheta) \,.
\end{equation*}
The first term follows from
$\nabla_{\vectheta}|\vectheta - \vecbeta|^2 = 2(\vectheta - \vecbeta)$.

\textbf{(b)} Setting $\nabla_{\vectheta}\,\tau = 0$:
\begin{equation*}
    (\vectheta - \vecbeta) - \nabla\psiL(\vectheta) = 0
    \qquad \Longrightarrow \qquad
    \vecbeta = \vectheta - \nabla\psiL(\vectheta) \,.
\end{equation*}
Since $\nabla\psiL(\vectheta) = \vecalpha(\vectheta)$ (the reduced
deflection angle), this is the lens equation
$\vecbeta = \vectheta - \vecalpha(\vectheta)$.  \hfill $\square$

\textbf{(c)} Physically, Fermat's principle states that light follows
paths of stationary travel time.  Among all possible positions
$\vectheta$ in the lens plane through which a photon could pass, only
those for which the arrival time is stationary (not changing to first
order under small displacements $\delta\vectheta$) correspond to actual
light paths.  These are the image positions.  The geometric delay
(longer path) and gravitational delay (Shapiro effect) compete, and
images form where their combined effect is stationary.

% ----- Exercise 6.2 -----
\subsection*{Exercise~\ref{ex:point_mass_arrival}: Arrival Times for a Point Mass Lens}

\textbf{(a)} With $\thetaE = 1''$ and $\beta = 0.5''$, the image
positions are:
\begin{align*}
    \theta_+ &= \frac{0.5 + \sqrt{0.25 + 4}}{2}
    = \frac{0.5 + 2.0616}{2} = 1.2808'' \,, \\
    \theta_- &= \frac{0.5 - 2.0616}{2} = -0.7808'' \,.
\end{align*}

\textbf{(b)} In units of $\thetaE$ ($x = \theta/\thetaE$, $y = 0.5$):
$x_+ = 1.2808$, $x_- = -0.7808$.
\begin{align*}
    \tau(x_+) &= \half(1.2808 - 0.5)^2 - \ln(1.2808) \\
    &= \half(0.7808)^2 - 0.2474
    = 0.3048 - 0.2474 = 0.0574 \,, \\
    \tau(x_-) &= \half(-0.7808 - 0.5)^2 - \ln(0.7808) \\
    &= \half(1.2808)^2 - (-0.2474)
    = 0.8202 + 0.2474 = 1.0676 \,.
\end{align*}

\textbf{(c)} The dimensionless time delay is:
\begin{equation*}
    \Delta\tau = \tau(x_-) - \tau(x_+) = 1.0676 - 0.0574 = 1.0101 \,.
\end{equation*}
This matches the analytic formula:
$\Delta\tau = y\sqrt{y^2+4}/2 + \ln[(\sqrt{y^2+4}+y)/(\sqrt{y^2+4}-y)]
= 0.5154 + 0.4947 = 1.0101$.  \checkmark

\textbf{(d)} For $M = 10^{12}\,\Msun$, $z_d = 0.3$:
\begin{align*}
    \RS &= \frac{2\Gnewton M}{\speedoflight^2}
    = \frac{2 \times 6.674 \times 10^{-11} \times
    1.989 \times 10^{42}}{(2.998 \times 10^8)^2}
    \approx 2.953 \times 10^{15}~\text{m} \,, \\
    \Delta t &= \frac{\RS}{\speedoflight}\,(1 + z_d)\,\Delta\tau
    = \frac{2.953 \times 10^{15}}{2.998 \times 10^8}
    \times 1.3 \times 1.0101 \\
    &\approx 1.293 \times 10^{7}~\text{s}
    \approx 149.7~\text{days} \,.
\end{align*}

% ----- Exercise 6.3 -----
\subsection*{Exercise~\ref{ex:typeI_first}: Type~I Image Arrives First}

\textbf{(a)} For a point mass lens, $\tau(x) = \half(x - y)^2 - \ln|x|$.
The image at $x_+$ is the primary image (minimum of $\tau$), and
$x_-$ is the secondary image (saddle point).  Since $x_+$ is a minimum
and $x_-$ is a saddle point, by definition
$\tau(x_+) \leq \tau(x_-)$ for any source position.

More explicitly: $\Delta\tau = y\sqrt{y^2+4}/2 +
\ln[(\sqrt{y^2+4}+y)/(\sqrt{y^2+4}-y)]$.  Both terms are positive for
$y > 0$: the first term is manifestly positive, and the argument of the
logarithm exceeds 1 since $\sqrt{y^2+4} + y > \sqrt{y^2+4} - y > 0$.
Therefore $\Delta\tau > 0$, i.e., $\tau(x_-) > \tau(x_+)$.
\hfill $\square$

\textbf{(b)} Numerical verification:

\begin{center}
\begin{tabular}{rcccc}
\toprule
$y = \beta/\thetaE$ & $\tau(x_+)$ & $\tau(x_-)$ & $\Delta\tau$ \\
\midrule
0.1  & $-0.0024$ & $0.0076$  & $0.0100$ \\
0.5  & $0.0574$  & $1.0676$  & $1.0101$ \\
1.0  & $0.1903$  & $2.6180$  & $2.4277$ \\
2.0  & $0.3508$  & $5.3573$  & $5.0065$ \\
\bottomrule
\end{tabular}
\end{center}

In all cases $\tau(x_+) < \tau(x_-)$.  \checkmark

\textbf{(c)} Since $t(\vectheta) \propto \tau(\vectheta)$ (with a
positive proportionality constant), a smaller $\tau$ means a smaller
arrival time.  Therefore the Type~I image (minimum of $\tau$) arrives
before the Type~II image (saddle point of $\tau$).  The time delay
between them is $\Delta t \propto \Delta\tau > 0$.

% ----- Exercise 6.4 -----
\subsection*{Exercise~\ref{ex:H0_from_time_delay}: Measuring $\Hubble$ from Time Delays}

\textbf{(a)} From the time delay formula:
\begin{equation*}
    \Delta t = \frac{(1 + z_d)}{\speedoflight}\,
    \frac{\Dd\,\Ds}{\Dds}\, (\tau_1 - \tau_2) \,.
\end{equation*}
Converting $\tau_1 - \tau_2 = 25~\text{arcsec}^2$ to radians:
$25 \times (1/206265)^2 = 5.876 \times 10^{-10}~\text{rad}^2$.
Converting $\Delta t = 420$~days $= 3.629 \times 10^7$~s.
With $z_d = 0.3$:
\begin{equation*}
    \frac{\Dd\,\Ds}{\Dds}
    = \frac{\Delta t \cdot \speedoflight}{(1 + z_d) \cdot (\tau_1 - \tau_2)}
    = \frac{3.629 \times 10^7 \times 2.998 \times 10^8}
    {1.3 \times 5.876 \times 10^{-10}}
    \approx 1.424 \times 10^{25}~\text{m} \,.
\end{equation*}
Converting to Mpc: $1.424 \times 10^{25} / 3.086 \times 10^{22}
\approx 4614$~Mpc.

\textbf{(b)} For a flat $\Lambda$CDM cosmology with $\Omegam = 0.3$, the
angular diameter distances scale as $D \propto \speedoflight / \Hubble$.
The dimensionless comoving distances (from Module~3) are $d_d(z_d=0.3)$,
$d_s(z_s=1.0)$, and $d_{ds}$.  Using numerical integration:
\begin{align*}
    d_d &\approx 0.291 \,, \quad d_s \approx 0.480 \,, \quad d_{ds} \approx 0.274 \,, \\
    \frac{\Dd\,\Ds}{\Dds} &= \frac{\speedoflight}{\Hubble}\,
    \frac{d_d\,d_s / [(1+z_d)(1+z_s)]}{d_{ds} / (1+z_s)}
    = \frac{\speedoflight}{\Hubble}\,
    \frac{d_d\,d_s}{(1+z_d)\,d_{ds}} \\
    &= \frac{\speedoflight}{\Hubble}\,
    \frac{0.291 \times 0.480}{1.3 \times 0.274}
    = \frac{\speedoflight}{\Hubble} \times 0.3923 \,.
\end{align*}

\textbf{(c)} Combining:
\begin{equation*}
    \Hubble = \frac{\speedoflight \times 0.3923}
    {\Dd\,\Ds/\Dds}
    = \frac{2.998 \times 10^8 \times 0.3923}{1.424 \times 10^{25}}
    \approx 8.25 \times 10^{-18}~\text{s}^{-1} \,.
\end{equation*}
Converting to km\,s$^{-1}$\,Mpc$^{-1}$: multiply by $3.086 \times 10^{22}$
(m per Mpc) and divide by $10^3$ (m per km):
\begin{equation*}
    \Hubble \approx 8.25 \times 10^{-18} \times 3.086 \times 10^{19}
    \approx 254.6~\text{km\,s}^{-1}\,\text{Mpc}^{-1} \,.
\end{equation*}

\begin{remark}
The high value reflects the simplified setup (the ``$\tau_1 - \tau_2$''
and distance model are illustrative).  In a more detailed treatment with
self-consistent lens modeling and precise distance ratios, the H0LiCOW
program obtains $\Hubble = 73.3^{+1.7}_{-1.8}$~km\,s$^{-1}$\,Mpc$^{-1}$,
which is consistent with local distance ladder measurements but in
tension with the Planck CMB value of $\Hubble \approx 67.4$~km\,s$^{-1}$\,Mpc$^{-1}$.
\end{remark}

% ----- Exercise 6.5 -----
\subsection*{Exercise~\ref{ex:arrival_surface_plot}: Plotting the Arrival-Time Surface}

\textbf{(a)--(b)} The arrival-time surface
$\tau(\theta_1, \theta_2)$ for a point mass lens with
$\vecbeta = (0.3, 0)\,\thetaE$ and $\thetaE = 1$ is:
\begin{equation*}
    \tau(\theta_1, \theta_2) = \half\left[(\theta_1 - 0.3)^2 + \theta_2^2\right]
    - \ln\sqrt{\theta_1^2 + \theta_2^2} \,.
\end{equation*}
The two stationary points are the image positions:
\begin{align*}
    \theta_+ &= \frac{0.3 + \sqrt{0.09 + 4}}{2} = 1.1612 \quad
    \text{(minimum of $\tau$)} \,, \\
    \theta_- &= \frac{0.3 - \sqrt{0.09 + 4}}{2} = -0.8612 \quad
    \text{(saddle point of $\tau$)} \,.
\end{align*}

\textbf{(c)} The contour plot shows nested contours around the minimum
at $(\theta_+, 0) = (1.1612, 0)$, and a saddle pattern at
$(\theta_-, 0) = (-0.8612, 0)$.  The Einstein ring
$|\vectheta| = \thetaE = 1$ passes between the two stationary points,
as expected (the primary image is outside $\thetaE$, the secondary
inside).

\textbf{(d)} These match the analytical predictions from the lens
equation $\theta_\pm = (\beta \pm \sqrt{\beta^2 + 4\thetaE^2})/2$
exactly.  \checkmark

\noindent
The plots are generated in
\mathlink{06_Fermat_Time_Delays/time_delay_function.wl}{time\_delay\_function.wl}
and the interactive version is in
\mathlink{06_Fermat_Time_Delays/arrival_time_surface.nb}{arrival\_time\_surface.nb}.
